
\documentclass[12pt,a4paper]{article}

% -----------------------------------------------------------------------------
% RSC / Digital Discovery submission-style source
% -----------------------------------------------------------------------------
% Conversion note: v16 replaced the ACS achemso wrapper used in earlier working
% drafts with an RSC-compatible article source using the rsc citation package.
% v18 confirms the MIT software license and prepares the DOI-ready release package while preserving the
% scientific content, tables, software-resource framing, and validated outputs.
% Final journal portal requirements should still be checked before submission.

\usepackage[utf8]{inputenc}
\usepackage[T1]{fontenc}
\usepackage[a4paper,margin=1in]{geometry}
\usepackage{amsmath}
\usepackage{amssymb}
\usepackage{booktabs}
\usepackage{graphicx}
\usepackage{newtxtext,newtxmath}
\usepackage[version=3]{mhchem}
\usepackage{url}
\usepackage{xcolor}
\usepackage[super]{rsc}
\usepackage[colorlinks=true,linkcolor=blue,citecolor=blue,urlcolor=blue]{hyperref}

\title{A Reproducible Software-Resource Workflow for AI-Guided Prediction of Monolithic Metal--Organic Framework Formation}
\author{Jinnan Wei\\
Department of Chemistry, University of Cambridge, Cambridge, UK\\
E-mail: jw2292@cam.ac.uk}
\date{}

\begin{document}
\maketitle

\begin{abstract}
Metal--organic frameworks (MOFs) are a versatile class of crystalline porous materials, yet most are isolated as microcrystalline powders with limited bulk density, mechanical robustness, and device-level processability. Monolithic MOFs (monoMOFs) address this morphological bottleneck by assembling nanoscale crystallites into binder-free, centimeter-scale architectures, but their discovery remains empirical because gelation, drying, capillary stress, and framework stability are coupled across multiple length and time scales. Here, we present a reproducible software workflow for AI-guided monoMOF discovery that combines literature text-mining interfaces, structured synthesis descriptors, gelation physics, and a Monolithic Formation Probability score (MFP-Score). MFP-Score rankings, OPR suggestions, feature-importance diagnostics, uncertainty classes, ablation summaries, sensitivity checks, output manifests, and checksums are generated by the accompanying synthetic-curated Supplementary Software~1 package. The framework defines how precursor concentration, solvent polarity, modulator acidity, drying kinetics, and optional physics-informed descriptors can be mapped into candidate Optimal Processing Routes (OPRs), while keeping all candidate routes hypothesis-generating rather than experimentally validated. This article is therefore positioned as a Digital Discovery-oriented software-resource manuscript: it provides a transparent prototype, one-command regeneration script, output validator, traceability map, and reproducibility interface, while reserving high-fidelity DFT/MD, autonomous-lab execution, and experimental route validation for future work. The work aims to convert monoMOF discovery from ad hoc empirical trial-and-error into a transparent, data-structured, and reproducible materials-informatics workflow.
\end{abstract}


% --- MAIN BODY ---

\section{Introduction: The Morphological Bottleneck in MOF Engineering}

Metal–organic frameworks (MOFs) are crystalline materials composed of inorganic nodes and organic linkers, forming extended porous networks with highly tunable functionalities.\cite{ref1} Their exceptional surface areas and chemical modularity enable applications in gas storage, catalysis, separation, and sensing. However, most MOFs are obtained as fine powders via solvothermal or hydrothermal synthesis, which severely limits industrial deployment. The morphological bottleneck constitutes a major challenge, particularly in large-scale processes: powdered MOFs exhibit high pressure drops in industrial columns, leading to energy inefficiencies, and their low mechanical strength causes attrition, generating fine dust that can obstruct equipment, contaminate systems, and pose operational and safety risks. Conventional densified pellets offer only a partial solution, as polymeric or graphite binders (typically 10–20 wt\%) reduce accessible porosity and volumetric capacity.

Monolithic MOFs (monoMOFs) address these limitations by assembling nanoscale crystallites into centimeter-scale, mechanically robust solids without binders or high-pressure compaction.\cite{ref2} These binder-free architectures introduce hierarchical porosity, combining intrinsic microporosity with inter-particle mesopores. Monolithic forms of Zr-based UiO-66\cite{ref4} and Cu-based HKUST-1\cite{ref5} have demonstrated record volumetric gas uptake. Importantly, monolithic morphology is not assumed to universally improve every property; rather, it shifts the relevant performance basis from gravimetric powder metrics to application-specific volumetric, mechanical, transport, and processability metrics. Nevertheless, the discovery rate of monoMOFs remains limited, reflecting a significant knowledge gap: the mechanisms governing the gel-to-monolith transition are poorly understood.\cite{ref6} To date, only a handful of pristine or composite monoMOF examples have been fabricated, including SnO$_2$ composites of ZIF-8\cite{ref7} and broader MOF gel/monolith families.\cite{ref6} Capturing the complex, multivariate dependencies of monoMOF formation—spanning precursor concentration, solvent polarity ($\pi^*$), modulator identity ($pK_a$), and drying kinetics—exceeds the capabilities of manual optimization.\cite{ref9}

Artificial intelligence (AI), particularly large language models (LLMs) and machine learning (ML), offers a practical route for organizing this multivariate synthesis space. Inspired by recent advances in AI-assisted MOF synthesis extraction and prediction,\cite{ref3,ref10,ref16} we hypothesize that monoMOF formation can be approached as a structured, testable ranking problem rather than as a purely empirical search. Here, we present a prototype-supported framework for AI-guided monoMOF discovery: (1) automated literature extraction via LLMs, (2) numerical feature encoding grounded in gelation physics, and (3) candidate Optimal Processing Route (OPR) generation for experimental prioritization. We introduce a quantitative descriptor--the MFP-Score--as a proposed probability-like ranking score for monolith formation from powder MOF precursors. MFP-Score values and OPRs are treated as hypothesis-generating outputs that require further curated-data training and experimental validation before they can be claimed as predictive performance metrics. The accompanying Supplementary Software package provides a minimal reproducible prototype for this interface and is intended to make the framework inspectable by reviewers rather than to claim a fully autonomous or experimentally validated materials-discovery system. For the present Digital Discovery-oriented version, the software package has been extended with \texttt{run\_all.py}, \texttt{validate\_outputs.py}, environment files, output manifests, checksums, ablation diagnostics, and sensitivity summaries so that the manuscript tables can be regenerated and checked from a clean software workflow.

\paragraph{Scope of the present article.} The present article is framed as a methods/software-resource computational workflow with a synthetic-curated demonstration package. The contribution is the definition of a reproducible MFP-Score/OPR interface, a one-command software workflow for regenerating manuscript-facing outputs, and the validation roadmap needed to convert the framework into a fully predictive monoMOF discovery platform. The framework should therefore be evaluated as a reproducible materials-informatics workflow, not as a conventional experimental synthesis report or as a completed high-throughput DFT/MD study.

\paragraph{Digital Discovery submission framing.} For a Digital Discovery route, the main manuscript is framed around reusable software, machine-readable outputs, transparent data/code availability, and a release-ready repository package rather than around claims of experimentally completed synthesis. The accompanying package therefore includes not only the MFP-Score and OPR scripts, but also a reviewer quick-start guide, environment files, output manifests, checksums, validation reports, and a manuscript-output traceability map. These components are intended to make the workflow inspectable, re-runnable, and prepared for deposition in a persistent repository with a versioned DOI before submission or acceptance.


\section{Theoretical and Mathematical Framework of MonoMOF Formation}

MonoMOF formation is fundamentally a two-stage kinetic process: (i) nucleation and growth of nanocrystals within a sol–gel medium (based on traditional sol-gel chemistry\cite{ref11} and non-silica monolithic synthesis\cite{ref12}), and (ii) slow solvent removal leading to nanoparticle coalescence into a macroscopic solid. The goal of monolithic synthesis is to maintain kinetic control, favoring a stable gel over rapid precipitation.

\subsection{Thermodynamic and Kinetic Perspectives}

The gelation process follows classical nucleation theory modified for confined growth, where the nucleation rate ($J$) is given by:
\[
J = A \cdot \exp\left(-\frac{\Delta G^*}{k_B T}\right)
\]
where $A$ is a kinetic prefactor, $k_B$ the Boltzmann constant, and $\Delta G^*$ the critical free energy barrier. For monoMOFs, the goal is to maintain $J$ at an intermediate regime—enabling abundant nanocrystal formation yet allowing interparticle connectivity. A crucial divergence in monoMOF synthesis is the transition from classical growth to non-classical aggregation. Unlike classical LaMer growth, where monomers add to a growing crystal, monolithic formation relies on \textit{Oriented Attachment (OA)} or \textit{Particle-Mediated Crystallization (PMC)}. In this regime, small, primary nanocrystals ($<10$ nm) are the building blocks. These nanocrystals collide and attach, driven by short-range forces (van der Waals, H-bonding) and guided by specific crystal facets, to form a larger, interconnected network.\cite{ref50} The kinetics of aggregation ($k_{agg}$) must be carefully balanced against the kinetics of dissolution and recrystallization ($k_{diss}$). If $k_{agg}$ is too fast, an unstable, flocculated precipitate (powder) forms. If $k_{diss}$ is dominant (e.g., in the presence of strong modulators or at high temperatures), Ostwald ripening occurs, where large crystals grow at the expense of smaller ones, again leading to a powder. The gel window is a specific kinetic regime where $k_{agg}$ is fast enough to form a sample-spanning network but slow enough to allow particles to align (OA), forming strong, crystalline `necks' between particles.

Gel densification is governed by capillary stress during drying. The capillary pressure ($P_c$) acting on the gel network, which, if too high, causes network collapse, is approximated by the Young-Laplace equation:
\[
P_c \approx \frac{2 \gamma \cos\theta}{r_p}
\]
where $\gamma$ is the solvent surface tension, $\theta$ the contact angle, and $r_p$ the pore radius. Lower $P_c$ (achieved via low-surface-tension solvents or supercritical fluid drying) prevents network collapse. The relationship between the particle-particle interaction potential ($V_{total}$) and the critical aggregation concentration ($CAC$) determines the gel's mechanical stability before drying. $V_{total}$ is modeled by the DLVO theory, combining electrostatic repulsion ($V_{elec}$) and van der Waals attraction ($V_{vdw}$):
\[
V_{total} = V_{elec} + V_{vdw}
\]
The proposed ML architecture is designed to learn the balance of these forces by weighting the solvent polarity ($\pi^*$) and modulator concentration ($\alpha$) once a reproducible training dataset is implemented.

\subsection{Advanced Mathematical Formulation and AI Integration}

We define a \textit{Monolithic Formation Function (MFF)}, $F_m$, incorporating key experimental variables:
\[
F_m = f(C_m, C_l, \eta, \phi, \alpha, T, R_d, \text{Solvent } \pi^*, \text{Modulator } pK_a, \Delta G_{solv}, \text{Node Stability})
\]
where $C_m$ and $C_l$ are metal and linker concentrations, $\eta$ solution viscosity, $\phi$ nanoparticle volume fraction, $\alpha$ modulator ratio, $T$ synthesis temperature, and $R_d$ solvent removal rate. New parameters include the Solvation Free Energy ($\Delta G_{solv}$) of the linker, which controls solubility, and a \textit{Node Stability Index} that could be derived in future versions from DFT calculations, curated literature descriptors, or experimentally inferred stability proxies. Empirically, monolith formation is favored when viscosity $\eta$ and volume fraction $\phi$ exceed critical values, and the removal rate $R_d$ is below a critical threshold $R_{crit}$. The eventual ML model is intended to learn the complex, non-linear dependencies within $F_m$ that determine $R_{crit}$ once a curated training set is available.

The probability of monolith formation $P_m$ scales sigmoidally with $F_m$. This logistic relationship defines the MFP-Score (ranging from 0 to 1):
\[
P_m = \frac{1}{1 + \exp(-(\beta_0 + \mathbf{\beta}^T \cdot \mathbf{x}))}
\]
where $\mathbf{x}$ is the multi-dimensional feature vector, and $\mathbf{\beta}$ represents the vector of learned coefficients, incorporating the complexity of the feature transformation $\phi(\mathbf{x})$.

Furthermore, the formation of a monolithic gel can be modeled by a time-dependent concentration field $C(r, t)$ of precursor species, governed by a reaction–diffusion equation:
\[
\frac{\partial C}{\partial t} = D \nabla^2 C - k_f C^n + k_b(C_g - C)^m
\]
where $D$ is the diffusion coefficient, $k_f$ the forward nucleation rate constant, $n$ the reaction order, and $k_b(C_g - C)^m$ a backward dissolution/re-aggregation term. This backward term is critical: it represents the Ostwald ripening process (where $C_g$ is the equilibrium concentration at a particle surface, dependent on curvature), which is the primary competitor to gelation. The intended AI task is to find conditions (modulators, solvents) that suppress this term, favoring the $k_f$ (nucleation) and $k_{agg}$ (aggregation) pathways. The concentration field $C(r, t)$ must be analyzed in the context of a \textit{critical depletion zone} around growing crystallites; if this zone is too large, it leads to inhomogeneous growth and powder formation. Numerical integration yields the gelation timescale $\tau_g$, which correlates inversely with the MFP-Score via the empirical fit, providing a route toward \textit{kinetic prediction} once calibrated:
\[
\tau_g \approx \tau_0 \cdot \exp\left(-\gamma \cdot P_m\right)
\]

\subsection{Gel Window and Experimental Variable Interpretation}

The kinetic regime that favors monoMOF formation, often referred to as the \textit{gel window}, is defined by a delicate balance between aggregation ($k_{agg}$) and dissolution/recrystallization ($k_{diss}$) rates. Within this window, $k_{agg}$ is sufficiently fast to establish a continuous network of primary nanocrystals, while $k_{diss}$ is slow enough to prevent Ostwald ripening from dominating. Deviations from this regime lead to either powder formation (if $k_{agg} \ll k_{diss}$) or inhomogeneous, fragile gels (if $k_{agg} \gg k_{diss}$).

The Monolithic Formation Function (MFF), $F_m$, quantitatively encodes the influence of key experimental parameters on the formation of a stable gel network. In practical terms: 

\begin{itemize}
    \item Metal and linker concentrations ($C_m$, $C_l$): Higher concentrations increase the collision frequency of nanocrystals, promoting network connectivity but can also accelerate undesired precipitation if excessive. 
    \item Modulator ratio ($\alpha$) and identity ($pK_a$): Modulators slow nucleation and growth, tuning particle size and interparticle attachment; strong modulators or high ratios favor small, uniform crystallites conducive to stable monoliths.
    \item Solvent polarity ($\pi^*$) and solvation energy ($\Delta G_{solv}$): These control solubility and particle–particle interactions; appropriate solvents stabilize the gel network while minimizing aggregation into powders.
    \item Solution viscosity ($\eta$) and nanoparticle volume fraction ($\phi$): Higher viscosity slows particle motion, reducing premature aggregation, while sufficient particle volume fraction ensures percolation and network formation.
    \item Temperature ($T$): Regulates both nucleation and aggregation kinetics; moderate temperatures maintain the gel window, whereas too high or low temperatures may induce uncontrolled precipitation or weak networks.
    \item Solvent removal rate ($R_d$): Controls capillary stresses during drying; slow removal preserves gel integrity, whereas rapid removal leads to network collapse and powder formation.
    \item Node stability (from future DFT/literature/experimental stability descriptors): Chemically stable nodes resist dissolution, supporting network connectivity during gelation and drying.
\end{itemize}

By integrating these variables, $F_m$ captures the multi-dimensional dependencies that determine the MFP-Score. High MFP-Scores should be interpreted as ranking indicators for conditions aligned with the gel window, not as validated probabilities of successful monoMOF formation until calibrated against experimental data.

\subsection{Thermodynamic Descriptor Space for AI Feature Learning}

The AI model maps synthesis parameters into a high-dimensional feature space $\mathbf{x}$, incorporating experimentally derived and calculated descriptors such as: Gelation energy ($\Delta G_{gel}$), Drying energy density ($E_{dry}$), Solvent polarity index ($\pi^*$), Surface tension ($\gamma$), Contact angle ($\theta$), and Metal–ligand coordination entropy term ($\Delta S_{coord}$). The feature space $\mathbf{x}$ is partitioned into three key categories:
\begin{enumerate}
    \item Chemical/Structural Features: Properties of the metal node (effective charge, coordination number, cluster size) and the linker (flexibility index, number of binding sites, GNN embeddings).
    \item Process/Kinetic Features: $C_m$, $C_l$, $\alpha$, $T$, $R_d$ (expressed in units of $\text{mm} \cdot \text{h}^{-1}$ for evaporation).
    \item Solvent/Environment Features: $\pi^*$, $\gamma$, $pK_a$ of the modulator, and the full multi-component solvent composition vector.
\end{enumerate}
A future implemented machine learning algorithm would learn a mapping $P_m = f(\mathbf{x})$,\cite{ref13} which can be expressed as a neural function approximation:
\[
P_m = \sigma \left( \mathbf{w}^T \phi(\mathbf{x}) + b \right)
\]
where $\phi(\mathbf{x})$ is a nonlinear kernel transformation (e.g., from a \textit{Graph Neural Network (GNN)} on the linker structure), $\mathbf{w}$ the trainable parameters, and $\sigma$ the logistic activation. The use of GNNs is vital as they can effectively encode the non-local connectivity and topological information of the organic linkers, which is a major determinant of the resulting MOF topology and mechanical resilience.\cite{ref14}

\subsection{Optional Physics-Informed Descriptor Interface: Quantum Chemistry and DFT}

To provide a future ML implementation with a deeper understanding of intrinsic chemical reactivity and stability, we define an optional physics-informed descriptor interface based on DFT-accessible quantities. These descriptors are not claimed here as completed high-throughput calculations; rather, they specify which electronic-structure features could be added in a later computational layer to move beyond purely empirical parameters.

\subsubsection{Node and Linker Stability Features}

The stability of the inorganic node ($\text{Zr}_6, \text{Fe}_3, \text{Zn}_4$ clusters) and the organic linker in the presence of solvent and modulator is paramount. A future DFT descriptor layer could provide:
\begin{itemize}
    \item Frontier Molecular Orbital (FMO) Analysis: The energy gap between the Highest Occupied Molecular Orbital ($\text{E}_{HOMO}$) and the Lowest Unoccupied Molecular Orbital ($\text{E}_{LUMO}$) of the organic linker provides an indirect measure of its chemical hardness and stability against side reactions. A smaller gap suggests higher reactivity.
    \[
    E_{gap, Linker} = E_{LUMO} - E_{HOMO}
    \]
    \item Nucleophilicity and Electrophilicity Indices: Based on the FMO energies, these indices ($v$ and $\omega$) quantify the linker's tendency to donate or accept electrons during coordination, directly influencing the kinetic rate constant $k_f$ in the reaction-diffusion equation (Section 2.2).
    \[
    \omega \approx \frac{\mu^2}{2\eta}, \text{ where } \mu \text{ is the chemical potential and } \eta \text{ the chemical hardness.}
    \]
    \item Cluster Dissociation Energy ($\Delta E_{diss}$): For the inorganic node (e.g., $\text{Zr}_6 \text{O}_4 (\text{OH})_4$ cluster), DFT could be used to estimate the energy required to remove one metal atom or cleave one bridging group. This energy serves as a crucial input for the Node Stability Index, potentially informing future rankings for high-temperature synthesis conditions $T$.
    \[
    \text{Node Stability Index} \propto 1 / \Delta E_{diss}
    \]
    \item Electrostatic Potential (ESP) Surfaces: Analysis of the ESP map of the linker molecule identifies the most favorable sites for protonation/deprotonation (influenced by the modulator $pK_a$) and subsequent coordination to the metal cluster. Descriptors include the minimum and maximum ESP values around the functional groups.
\end{itemize}

\subsubsection{Solvation and Intermolecular Interaction Energies}

The transformation from a stable sol to a monolithic gel is primarily dictated by subtle variations in particle–particle and particle–solvent interactions. a future Density Functional Theory (DFT) layer could provide a rigorous framework to quantify these effects at the electronic level:

1. Solvation Free Energy ($\Delta G_{solv}$)  

   Implicit solvation models such as COSMO or PCM could be used to estimate $\Delta G_{solv}$ for both the linker and the metal precursor within the predicted solvent environment. Such calculations would capture the thermodynamic balance between solvation and precipitation, directly informing the optimal precursor concentrations ($C_m$ and $C_l$). By including solvent dielectric properties, such computations would also account for differential stabilization of charged or polar intermediates, which can influence nucleation and early-stage aggregation.

2. Hydrogen Bonding and Intermolecular Interactions  

   For modulators and protic solvents (e.g., water, formic acid), future DFT calculations on small molecule–cluster complexes could quantify hydrogen bonding strengths and non-covalent interactions (e.g., $\pi$–stacking, van der Waals contacts). These interactions govern the kinetics of particle association ($k_{agg}$) and the propensity for gel network formation. Evaluating the binding energies and optimized geometries would provide insight into which solvent–modulator combinations favor controlled aggregation and monolith stability.

By integrating these quantum chemical descriptors in a future implementation, the model could capture electronic-level drivers of monoMOF formation more explicitly. This multi-scale approach—linking molecular energetics to macroscopic morphology descriptors such as the MFP-Score—could enhance the mechanistic interpretability and predictive power of the machine learning framework, moving beyond purely correlative modeling toward chemically informed prediction.


\subsection{Future High-Fidelity Layer: TST and Molecular Dynamics}

A later high-fidelity implementation could provide more predictive, rather than purely correlative, features by adding Transition State Theory (TST) for nucleation kinetics and Molecular Dynamics (MD) for mechanical stability. In the present manuscript, these methods are presented as a future descriptor layer and computational roadmap, not as completed calculations supporting the current MFP-Score values. Such calculations would require explicit input structures, software versions, force-field parameterization, convergence criteria, and archived outputs before they could be treated as quantitative results.

\subsubsection{Transition State Theory (TST) for Nucleation Kinetics}

The nucleation rate prefactor $A$ and activation barrier $\Delta G^*$ (Section 2.1) are highly sensitive to the initial coordination between the metal precursor and the linker, which can be systematically quantified using TST. This approach enables a physically grounded estimation of the reaction kinetics, rather than relying on empirical correlations. The scenario parallels zeolite synthesis, where organic structure-directing agents (OSDAs) play a role analogous to MOF modulators.\cite{ref15}

DFT Functional and Basis Set Selection:  
Accurate modeling of transition metal clusters (e.g., $\text{Zr}_6, \text{Fe}_3$) would require careful choice of density functionals and basis sets to capture both electronic correlation and dispersion interactions. Standard GGA functionals (PBE, BLYP) often underestimate weak non-covalent interactions; therefore, a future implementation could benchmark hybrid or dispersion-corrected functionals such as B3LYP-D3(BJ), meta-GGA functionals (M06-L, SCAN), or range-separated hybrids ($\omega\text{B97X-D}$) for representative transition states. Basis sets and pseudopotentials would need to be selected explicitly, with convergence criteria, solvent models, input structures, and output files archived before any DFT-derived descriptor is treated as a quantitative manuscript result.

Transition State Search:  
A future TST workflow would identify the geometry of the highest-energy point along the reaction coordinate (Transition State, TS) for the initial metal-linker bond formation. This computationally rigorous approach would allow quantification of the energy barrier that governs nucleation kinetics. The procedure involves:
\begin{enumerate}
    \item Reactants and Products Geometry Optimization: Ensuring that starting and ending structures are energetically minimized.
    \item Nudged Elastic Band (NEB) Method: Generating a chain of intermediate structures connecting reactants to products, approximating the reaction pathway.
    \item Dimer Method/Vibrational Analysis: Refining the TS and confirming it is a first-order saddle point via the Hessian matrix (one imaginary frequency corresponding to the reaction coordinate).
\end{enumerate}

The forward rate constant $k_f$ is calculated as:
\[
k_f = \kappa \frac{k_B T}{h} \exp\left(-\frac{\Delta G^\ddagger}{R T}\right)
\]
where $\Delta G^\ddagger$ includes zero-point energy and thermal corrections, $h$ is Planck’s constant, and $\kappa$ is the transmission coefficient (typically 1). This physically derived $k_f$ could then be used as a predictive feature in the reaction-diffusion ML model (Section 2.2), enabling a more physics-informed simulation of nucleation rates and synthesis conditions once the calculation is actually performed. Linking archived first-principles calculations with machine learning would help anchor future AI predictions in chemical physics rather than statistical correlation alone.


\subsubsection{Molecular Dynamics (MD) for Mechanical and Viscoelastic Properties}

The mechanical integrity of monoMOFs is critical, and future MD simulations could provide quantitative features related to the material's elastic response ($E_{comp}$) and thermal stability, which are challenging to obtain experimentally. Such simulations would complement the ML model by translating atomistic interactions into macroscopic mechanical descriptors.

Force Field Selection:  
MOFs are complex periodic structures with both inorganic clusters and organic linkers. Classical force fields (FFs) could be employed to model intra- and inter-molecular interactions efficiently. The Universal Force Field (UFF) provides broad coverage, but for higher accuracy in capturing flexibility and non-bonded interactions, DREIDING or custom reactive force fields derived from DFT (ReaxFF) are preferred. MOF-specific FFs such as MOF-FF offer the best accuracy for well-characterized systems like UiO-66 and ZIF-8.\cite{ref51}  
From a computational chemistry perspective, proper parameterization of the metal-linker coordination bond is essential because this interaction governs the mechanical response and defect formation. While harmonic potentials capture small oscillations, reactive force fields (ReaxFF) are required to model bond-breaking and bond-forming events, which influence long-term stability, failure modes, and gelation dynamics.

Simulating Elastic Properties:  
A future MD workflow would compute the elastic stiffness tensor ($C_{ijkl}$) by applying small orthogonal strains ($\epsilon_k$) to the unit cell and calculating the resulting stress ($\sigma_i$) using the virial theorem under NPT (isothermal–isobaric) or NVE (microcanonical) ensembles:
\[
\sigma_i = C_{ik} \epsilon_k
\]
This tensor would encode the directional response of the crystal lattice to applied strain. From it, aggregate mechanical properties are derived:
\begin{itemize}
    \item Bulk Modulus ($K$): Resistance to volumetric deformation.
    \item Shear Modulus ($G$): Resistance to shape change under shear.
    \item Young's Modulus ($E_{comp}$): Stiffness under uniaxial compression.
\end{itemize}
These ideal crystal properties would serve as input features for the ML model, which then learns the \textit{Morphology Degradation Factor} (MFD) describing the reduction in stiffness when an ideal crystal assembles into a monolithic structure:
\[
E_{comp, mono} = \text{MFD} \cdot E_{comp, crystal}
\]
The MFD is influenced by solvent polarity ($\pi^*$) and drying rate ($R_d$), reflecting the critical role of inter-particle necks and porosity. This parallels the Hall-Petch effect in nanocrystalline metals but is dominated here by nanoscale connectivity rather than grain boundaries.

Simulating Gel Collapse:  
Capillary forces during drying can deform or fracture the monoMOF network. future MD simulations could model two nanocrystals separated by a solvent bridge to estimate the potential energy as a function of inter-particle separation. This would allow calculation of the maximum tension the network withstands before collapse, providing a \textit{Capillary Collapse Threshold} ($P_{collapse}$). This is directly compared to the predicted capillary pressure $P_c$ from the Young-Laplace equation to refine drying protocols in the AI-generated OPR.

Coarse-Grained (CG) MD:  
To connect atomistic simulations with macroscopic gelation, CG-MD could represent each MOF nanoparticle as a single bead and the solvent implicitly or as larger beads. Interaction potentials could later be parameterized from all-atom MD or literature force-field data to approximate inter-particle forces, including DLVO-type electrostatic and van der Waals interactions. This would allow simulation of thousands of nanoparticles over microseconds, capturing network formation, fractal dimension, and shear response, which inform the rheology model (Section 5.3).\cite{ref52}

By integrating these computational chemistry insights—bonding energetics, elastic response, and capillary stability—future MD simulations could provide mechanistically grounded features that enhance the ML model's predictive power beyond pattern recognition.


\section{AI-Driven Workflow and Predictive Modeling}

\subsection{Literature Mining and Feature Engineering via LLMs: The ChemPrompt Methodology}

The workflow begins with rigorous literature curation, which constitutes the most time-intensive and error-prone step in conventional materials informatics. To address this, the proposed workflow uses advanced natural language processing (NLP) techniques based on prompt-engineered LLMs, following the \textit{ChemPrompt methodology}\cite{ref3}. Unlike traditional keyword-based extraction, this approach enables sequence-aware extraction of synthesis procedures, capturing the ordered dependencies among reagents, solvents, modulators, and reaction conditions, even when reported in heterogeneous natural language.  

1. LLM Architecture and Prompt Engineering  

   A future implementation could use a high-capacity LLM, tuned or prompted for chemistry literature, to parse complex MOF synthesis procedures. The ChemPrompt methodology integrates chain-of-thought (CoT) reasoning with a structured output schema, allowing the decomposition of ambiguous textual descriptions into machine-readable steps while minimizing errors. For example, the sentence "The precursors were dissolved in DMF and heated" is parsed to identify:  
   - chemical entities (precursors, DMF)  
   - functional roles (precursors = reactants; DMF = solvent)  
   - process parameters (e.g., heating; temperature unspecified → $T = \text{Not Specified}$)  
   - sequential order of operations (dissolution followed by heating)  

2. Prompt Structure  

   The prompt enforces systematic reasoning and structured output:  
   2.1 System Instruction: Define the LLM’s persona as a World-Class Materials Chemist and Data Scientist.  
   2.2 Task Definition: Extract complete synthetic sequences, including reactants, stoichiometry, solvents, reaction temperature, duration, and post-synthetic modifications.  
   2.3 Constraint Enforcement for monoMOFs: Highlight descriptors relevant to monolithic formation (e.g., gel, viscous suspension, slow evaporation, SCF drying). Explicit mention of a gel labels the sample as a Positive Monolith Candidate.  
   2.4 Schema Output: Enforce a strict JSON schema for ML pipeline integration, including nested objects for multi-step operations.

3. Example Prompt Logic  

   The LLM extraction schema is designed to identify the metal source ($\text{M}_{src}$), linker ($\text{L}_{src}$), molar ratio ($\alpha_{M:L}$), reaction time ($t_{rxn}$), and post-synthesis treatment ($\text{PST}$), outputting them as a structured JSON object: Metal, Linker, Stoichiometry, Temperature, Drying Method, and Morphology Indicator.

4. Data Normalization and Feature Encoding  

   The structured LLM output is converted into quantitative features for ML:  
   4.1 Chemical Identity Encoding: Map metal nodes to canonical clusters or oxidation states; convert organic linkers to SMILES and generate 256-dimensional Graph Embedding Vectors capturing topology, electronics, and sterics.  
   4.2 Stoichiometry Normalization: Standardize all molar concentrations relative to the metal center.  
   4.3 Process Variable Quantification: Map qualitative descriptors like "slow evaporation" to physically meaningful numerical values (e.g., solvent removal rate $R_d$), encode SCF drying as a binary feature with associated capillary stress mitigation.

This numeric, hierarchical style mirrors the structure used in the TST and MD sections, ensuring uniform presentation throughout the methods-style sections.

\subsection{Machine Learning Architecture and Optimization for Imbalanced Data}

The high-quality feature vector $\mathbf{x}$ is now ready for the ML pipeline. The primary architectural challenge stems from the extreme data imbalance, with only slightly more than a dozen positive samples (monoMOFs) compared to over 100,000 negative samples (powder MOFs). 

\subsubsection{Hierarchical Feature Integration and GNNs}

The ML architecture is designed to handle hierarchical features:
\begin{enumerate}
    \item Structural Features (GNN): A future supervised implementation could use a \textit{Message Passing Neural Network (MPNN)} to generate the feature embedding $\phi(\mathbf{x}_{linker})$ from the organic linker graph. The MPNN iteratively updates the feature vector of each atom based on its neighbors, providing a comprehensive representation of the linker's topology, electronic state, and functional group distribution. This is a common strategy in nanomaterials synthesis prediction.\cite{ref17}
    \item Process and Quantum Features: The 256-D GNN vector is concatenated with the 30-D vector of process variables (concentration ratios, $T$, $R_d$) and the 10-D vector of quantum-chemical descriptors ($E_{gap}, \Delta E_{diss}, \text{ESP}_{min}$). This combined, high-dimensional vector $\mathbf{x}_{final}$ serves as the input to the classifier.
\end{enumerate}

\subsubsection{Weighted Ensemble Classification: Gradient Boosted Trees}

Given the imbalance, a simple Neural Network would be heavily biased towards predicting 'powder' (MFP-Score near 0). A future supervised implementation could use a \textit{Gradient Boosted Tree (GBT) model}, such as LightGBM, as a final classifier after the current prototype establishes the feature table and baseline ranking interface. A GBT-style model is attractive for this task for several reasons: (1) Its interpretability (via SHAP) is crucial for building trust with chemists. (2) It naturally handles heterogeneous features (i.e., high-dimensional GNN embeddings, continuous scalar temperatures, and categorical solvent types) without extensive feature scaling, unlike NNs. (3) It is highly robust to noisy features, which are common in data extracted from literature.

The standard binary cross-entropy loss function $L$ (Section 2.2) is modified using \textit{Class-Weighting (CW)}, where the weight $w_{i}$ for a monolithic sample ($y_i=1$) is set equal to the inverse ratio of the class frequencies ($\approx 100,000 / 12 \approx 8333$).
\[
L = -\frac{1}{N} \sum_{i=1}^N w_i \left[y_i \log(P_{m,i}) + (1-y_i) \log(1-P_{m,i})\right]
\]
where $w_i = W_{mono}$ for $y_i=1$ and $w_i = W_{powder}$ for $y_i=0$, and $W_{mono} \gg W_{powder}$. This forces the model to heavily penalize False Negatives (predicting a monoMOF is a powder), which is the most detrimental error in this discovery phase. Class weighting is preferred over over-sampling techniques in the planned supervised implementation like \textit{SMOTE (Synthetic Minority Over-sampling Technique)} because SMOTE generates "synthetic" data points by interpolating between existing positive samples. In a high-dimensional, sparse chemical space, this interpolation can easily create chemically nonsensical or physically impossible synthesis protocols, poisoning the training set. CW would avoid this by increasing the importance of real, experimentally verified data without fabricating new chemical protocols.

\subsubsection{Hyperparameter Optimization and Interpretability}

Model performance should be maximized by optimizing hyperparameters (e.g., number of trees, learning rate, minimum samples per leaf) using Tree-structured Parzen Estimator (TPE)-based Bayesian optimization and evaluated with imbalance-aware metrics such as AUPRC, balanced accuracy, recall at fixed precision, and calibrated probability error. In the current manuscript, numerical performance values are reported only as synthetic demonstration diagnostics unless regenerated from a curated dataset and validated in a future study built on curated experimental data. The minimum acceptable implementation should report the data split, positive/negative label definitions, class weights or resampling strategy, cross-validation protocol, and uncertainty/calibration diagnostics.

Post-training interpretability should be used to test whether the model captures chemically plausible trends. A future SHAP or permutation-importance analysis should specifically test whether modulator identity/concentration, drying rate, solvent polarity, and capillary-stress descriptors dominate the ranking in a way consistent with gelation chemistry. The accompanying prototype package generates outputs that are integrated into the manuscript: \texttt{mfp\_scores.csv} ranks 12 candidate powder MOFs; \texttt{candidate\_opr\_table.csv} provides the corresponding hypothesis-generating processing windows; \texttt{feature\_importance.csv} reports weight-based diagnostic priorities, with the three largest demonstration weights assigned to drying\_route\_score (0.189), framework\_stability\_z (0.168), capillary\_safety (0.161); and \texttt{uncertainty\_report.csv} keeps all final accuracy, SHAP, OPR, and synthesis-success claims disabled until curated data and experimental validation are available. The synthetic calibration file reports demonstration-only metrics on 18 seed records (precision 0.588, recall 1.0, F1 0.741, balanced accuracy 0.562), which are retained as prototype sanity-check diagnostics rather than as manuscript performance claims.
The model performs two key tasks:
\begin{enumerate}
    \item Classification: Predicting the MFP-Score ($P_m$).
    \item Optimization: Generating the Optimal Processing Route (OPR), a prescriptive, quantitative synthesis protocol. This is achieved by taking the existing feature vector of a powder MOF ($\mathbf{x}_{powder}$) and performing a constrained optimization search within the feature space to find the optimal vector $\mathbf{x}_{OPR}$ that maximizes $P_m$, ensuring that the suggested changes (e.g., increasing modulator concentration, lowering $R_d$) are chemically and physically plausible.
\end{enumerate}

\subsection{Active Learning and Closed-Loop Refinement: Scientific Reinforcement Learning}

Future AI-generated synthesis suggestions should undergo human expert verification for chemical plausibility, avoiding unsafe or unrealistic conditions. This hybrid human-AI workflow could then be implemented in an \textit{Active Learning (AL)} or \textit{Scientific Reinforcement Learning (RL)} loop. AL is necessary because the cost of obtaining a new, high-quality, labeled data point (i.e., synthesizing a new MOF and characterizing it) is extremely high. The AI system's intended role is to suggest the next most \textit{informative} experiment, not just the highest probability success.\cite{ref18}

\subsubsection{Active Learning Policy: Uncertainty Sampling}

The AL policy $\mathcal{A}_{AL}$ selects the next experiment $\mathbf{x}_{next}$ based on the model's prediction uncertainty:
\[
\mathbf{x}_{next} = \arg \max_{\mathbf{x} \in \mathcal{U}} H(P_m(\mathbf{x}))
\]
where $\mathcal{U}$ is the set of unverified powder MOFs, and $H(P_m)$ is the Shannon entropy (uncertainty) of the prediction $P_m$:
\[
H(P_m) = -P_m \log_2(P_m) - (1-P_m) \log_2(1-P_m)
\]
This policy prioritizes MOFs whose MFP-Score is close to 0.5 (where the model is least certain whether they will form a monolith or a powder), ensuring that the limited experimental resources are directed toward regions of the feature space where data acquisition will yield the greatest information gain, thus rapidly shrinking the model's confidence interval.

\subsubsection{Reinforcement Learning Formulation}

For full autonomous control, the process is formalized as a \textit{Markov Decision Process (MDP)}, where the goal is to optimize the control policy $\pi^*$—mapping experimental state $s$ to action $a$—to maximize expected monolith yield:
\[
\pi^*(s) = \arg \max_a E\left[\sum_{t=0}^\infty \gamma^t R(s_t, a_t)\right]
\]
\begin{itemize}
    \item State Space ($s$): Defined by the complete set of input features $\mathbf{x}_{final}$ plus the current MOF morphology (e.g., $s_t = \{\mathbf{x}_{final, t}, \text{Morphology}_t\}$).
    \item Action Space ($a$): The set of prescriptive changes the robotic system can execute, defined as vectors of perturbations to the OPR (e.g., $\Delta C_m, \Delta T, \Delta R_d, \dots$).
    \item Reward Function ($R$): $R(s_t, a_t)$ is $+1$ for a successful monolithic synthesis (MFP-Score $> 0.75$) and $-1$ for a failed powder synthesis. Intermediate rewards are given based on key validation metrics (e.g., $+0.5$ for high $E_{comp}$ but low $V_{micro}$).
    \item Discount Factor ($\gamma$): Encourages the model to find efficient, fast synthesis routes ($\text{minimize } T_{synth}$).
\end{itemize}
The system would evolve dynamically toward improved predictive accuracy only after a real experimental feedback loop is established. In the current manuscript, policy-gradient reinforcement learning and self-driving-lab execution are presented as future integration points rather than as completed components of the MFP-Score workflow.\cite{ref19}

\subsection{Causal Inference and Explainable AI (XAI): Moving Beyond Correlation}

While a GBT-style model could provide an interpretable baseline, its reliance on correlative features limits its ability to explain why a monolith forms. A later version should therefore test a causal-inference layer using Bayesian networks or structural causal models; in the present manuscript, this layer remains a methodological extension rather than a completed causal discovery result.\cite{ref20}

\subsubsection{Causal Graphical Models and Interventions (Do-Calculus)}

We construct a \textit{Directed Acyclic Graph (DAG)} where nodes represent features ($\text{Modulator } pK_a$, $C_m$, $T$, $\pi^*$) and edges represent hypothesized causal links. For example, $\text{Modulator } pK_a \rightarrow \text{Crystallite Size} \rightarrow \text{Network Porosity} \rightarrow \text{MFP-Score}$. This graph is a \textit{Structural Causal Model (SCM)} that encodes our physical understanding of the system.

The causal model is used to answer "what if" questions using Judea Pearl's \textit{do-calculus}, which formalizes the concept of intervention. For instance: \textit{What is the predicted MFP-Score if we \textbf{intervene} to set the modulator ratio} $\alpha$ \textit{to 30, regardless of the initial synthesis conditions?} This is denoted $P(P_m | \text{do}(\alpha=30))$, which is fundamentally different from the standard conditional probability $P(P_m | \alpha=30)$. This is vital because a simple correlation might show high $\alpha$ and high $P_m$ are correlated, but only because both are caused by a third, confounding variable (e.g., a specific metal node). The do-calculus allows us to isolate the true causal effect of $\alpha$. CI allows the model to predict the outcome of a novel experiment that has never been performed before, as long as the intervention is structurally sound.

Implementation with Bayesian Networks:
In a future causal implementation, the causal graph could be trained as a \textit{Bayesian Network (BN)}, where the joint probability distribution over all variables can be factorized according to the graph structure:
\[
P(\mathbf{x}, P_m) = \prod_{i} P(X_i | \text{Parents}(X_i))
\]
After sufficient data are available, BN parameters could be estimated to simulate the effect of interventions and calculate the \textit{Average Causal Effect (ACE)} of a variable change (e.g., the average change in MFP-Score caused by a 5-unit decrease in solvent $\pi^*$). This would test whether OPR recommendations reflect plausible causal mechanisms rather than spurious correlations.

\subsubsection{SHAP for Feature Interaction and Transparency}

While future SHAP analysis could provide feature importance, a deeper XAI approach is needed to understand \textit{interactions}. a later implementation could use \textit{SHAP interaction values} to quantify how two features jointly contribute to the prediction. For example, the interaction value $\text{SHAP}_{Interaction}(\text{Modulator } pK_a, T)$ might be highly positive, indicating that the impact of a high $pK_a$ (weak modulator) is only effective when coupled with a low synthesis temperature $T$. This interaction is physically meaningful: weak modulators require less thermal energy to overcome the activation barrier for coordination.

\subsubsection{Counterfactual Reasoning for Chemical Insight}

Beyond simple interventions, the structural causal model (SCM) enables \textit{counterfactual reasoning}, the most advanced level of causal inference. This capability allows us to pose specific retrospective questions. For instance, for the successful synthesis of mono-UiO-66 (a real data point), one can ask: "What would the MFP score have been \textit{if} acetic acid had been used instead of formic acid, \textit{given that} all other parameters (temperature, solvent, etc.) remained exactly as in the original experiment?" This query, expressed as $P(P_m^{'} \mid \mathbf{x}^{'}, \text{do}(\text{modulator}=\text{acetic acid}))$, would enable a future causal version of the model to test a single altered parameter in silico, isolating candidate factors responsible for success or failure and providing chemists with more specific, experimentally testable hypotheses.

By providing causal and interaction diagnostics, the AI system would become a transparent decision-support tool that allows human chemists to validate the reasoning and refine the underlying physical models.

\subsection{Advanced Graph Neural Networks (GNNs) for Linker and Topology Encoding}

The initial GNN (MPNN) provides a good embedding, but MOF structural complexity demands more sophisticated graph architectures to accurately predict topology and resilience.

\subsubsection{Graph Attention Networks (GAT)}

The \textit{Graph Attention Network (GAT)} enhances the MPNN approach by weighting the messages passed between neighboring atoms/functional groups. Instead of uniformly averaging features from neighbors, the GAT applies an \textit{attention mechanism} to learn the relative importance of different neighbors based on their current features.

The attention coefficient $\alpha_{ij}$ between atom $i$ and its neighbor $j$ is calculated as:
\[
\alpha_{ij} = \frac{\exp(\text{LeakyReLU}(\mathbf{a}^T [\mathbf{W}\mathbf{h}_i || \mathbf{W}\mathbf{h}_j]))}{\sum_{k \in \mathcal{N}_i} \exp(\text{LeakyReLU}(\mathbf{a}^T [\mathbf{W}\mathbf{h}_i || \mathbf{W}\mathbf{h}_k]))}
\]
where $\mathbf{h}_i$ and $\mathbf{h}_j$ are the feature vectors, $\mathbf{W}$ is a shared weight matrix, and $\mathbf{a}$ is a shared attention vector.

Benefit: In MOFs, this allows the GAT to assign greater importance (attention) to the coordination groups of linkers (e.g., carboxylates, pyrazoles) while assigning lower importance to inert alkyl groups. For the monolithic formation task, a future GAT could be trained to focus on features that influence flexibility, such as rotatable bonds and linker length, which directly affect the porosity and mechanical stability of the resulting framework.

\subsubsection{Advanced Architectures (SchNet, DimeNet)}

While GAT captures connectivity, other architectures are superior for modeling 3D atomic environments, which are crucial for predicting linker-solvent and linker-cluster interactions.
\begin{itemize}
    \item SchNet: A continuous-filter convolutional GNN that operates on 3D coordinates, making it ideal for learning representations of linker conformation and flexibility.\cite{ref21}
    \item DimeNet: An equivariant GNN that explicitly encodes angles and dihedral angles, providing a rich representation of the linker's 3D geometry and stereoelectronic properties.\cite{ref22}
\end{itemize}
By replacing the simple MPNN with a pre-trained SchNet or DimeNet in a future implementation, the model could gain a deeper representation of the steric and electronic factors that govern $\Delta G_{solv}$ and $k_{agg}$, potentially improving future MFP-Score predictions after validation.

\subsubsection{Self-Supervised Learning (SSL) for Representation Pre-Training}

Given the extreme imbalance and the limited number of reported monoMOF samples, the capacity of any supervised model to learn a rich representation is inherently constrained. To mitigate this, a future implementation could employ \textit{Self-Supervised Learning (SSL)} on the extensive, readily available, unlabeled dataset of all known MOF structures, such as the over 100,000 entries in the Cambridge Structural Database (CSD).

\textbf{SSL Task: Contrastive Learning.} 
The GNN could be pre-trained on a self-supervised task such as \textit{Contrastive Learning (CL)}. For a given linker $\mathbf{x}_{\text{linker}}$, two augmented views ($\mathbf{x}'$, $\mathbf{x}''$) are generated, for example by randomly masking atoms or perturbing bond angles. The objective is to pull the representations of these two views (the "positive pair") closer together in the embedding space, while pushing them away from the representations of all other linkers in the mini-batch (the "negative pairs"). The loss function, typically the $\text{NT-Xent}$ loss, is defined as
\[
\mathcal{L}_i = -\log \frac{\exp(\text{sim}(\mathbf{z}_i', \mathbf{z}_i'') / \tau)}
{\sum_{j=1}^{2N} 1_{[j \neq i]} \exp(\text{sim}(\mathbf{z}_i, \mathbf{z}_j) / \tau)},
\]
where $\mathbf{z}$ denotes the projected embeddings, $\text{sim}$ represents cosine similarity, $\tau$ is the temperature parameter, and $1_{[j \neq i]}$ is the indicator function that excludes the anchor itself. After pre-training, the learned GNN weights would be transferred and fine-tuned on the small, labeled monoMOF dataset. This procedure is expected to improve the quality of the linker embeddings and may support more robust MFP-Score predictions once sufficient labeled monoMOF data are available.




\section{Case Studies: Prototype-Generated monoMOF Conversion Routes}

In the current manuscript, the case-study tables are no longer purely illustrative placeholders. They have been regenerated from the accompanying Supplementary Software~1 outputs using the synthetic-curated candidate table, the baseline MFP-Score script, the capillary-stress calculator, and the OPR ranking script. The resulting rankings remain hypothesis-generating and are not experimental validation data, but they are now internally traceable to \texttt{mfp\_scores.csv}, \texttt{candidate\_opr\_table.csv}, \texttt{feature\_importance.csv}, and \texttt{uncertainty\_report.csv}. Accordingly, Tables~\ref{tab:mfp-ranking} and \ref{tab:opr-summary} provide the manuscript-facing summary of the prototype-generated outputs, while the full machine-readable outputs remain in the supplementary software package.


% --- Table 1: prototype-generated candidate ranking ---
\begin{table}[htbp]
  \centering
  \caption{Prototype-generated MFP-Score candidate ranking from Supplementary Software~1. Values are generated from the synthetic-curated demonstration table \texttt{candidate\_powder\_mofs.csv} and the transparent baseline script \texttt{mfp\_score\_baseline.py}. They are suitable for prioritization and manuscript-method integration, but they are not validated probabilities of experimental monoMOF formation.}
  \label{tab:mfp-ranking}
  \scriptsize
  \resizebox{\textwidth}{!}{%
  \begin{tabular}{@{}llllllll@{}}
    \toprule
    Candidate & Metal node & Linker family & Topology & MFP-Score & Priority & $P_c$ / MPa & Drying risk \\
    \midrule
    NU-1000 & Zr6 & TBAPy & csq & 0.974 & High & 1.090 & low \\
    UiO-67 & Zr6 & BPDC & fcu & 0.970 & High & 1.576 & low \\
    MOF-808 & Zr6 & BTC & spn & 0.967 & High & 1.442 & low \\
    MOF-801 & Zr6 & fumarate & fcu & 0.956 & High & 2.247 & low \\
    PCN-222 & Zr6 & TCPP & csq & 0.950 & High & 1.197 & low \\
    HKUST-1 & Cu paddlewheel & BTC & tbo & 0.887 & High & 1.367 & low \\
    MIL-100(Fe) & Fe trimer & BTC & mtm & 0.879 & High & 1.768 & low \\
    ZIF-8 & Zn imidazolate & 2-mIm & sod & 0.872 & High & 1.851 & low \\
    Mg-MOF-74 & Mg chain & DOBDC & rod & 0.830 & High & 4.400 & medium \\
    MIL-101(Cr) & Cr trimer & BDC & mtm & 0.732 & High & 2.103 & low \\
    MOF-5 & Zn4O & BDC & pcu & 0.597 & Medium & 2.536 & low \\
    CAU-10 & Al chain & isophthalate & chain & 0.500 & Medium & 12.000 & high \\
    \bottomrule
  \end{tabular}%
  }
\end{table}




% --- Table 2: prototype-generated OPR summary ---
\begin{table}[htbp]
  \centering
  \caption{Prototype-generated OPR summary from Supplementary Software~1. The routes are hypothesis-generating outputs produced by \texttt{candidate\_ranking.py} and \texttt{opr\_optimizer.py}; compact wording is used here for readability, and the full generated strings remain in \texttt{candidate\_opr\_table.csv}.}
  \label{tab:opr-summary}
  \scriptsize
  \begin{tabular}{@{}p{1.45cm}p{0.9cm}p{2.9cm}p{1.8cm}p{3.8cm}p{1.1cm}p{1.1cm}@{}}
    \toprule
    Candidate & Score & Modulator screen & Temp. & Recommended drying route & Risk & Unc. \\
    \midrule
    NU-1000 & 0.974 & FA/AcOH, 1.5--3.0 M & 80--120 $^\circ$C & scCO$_2$ reference route & low & low \\
    UiO-67 & 0.970 & FA/AcOH, 1.5--3.0 M & 80--120 $^\circ$C & solvent exchange + slow-drying screen & low & low \\
    MOF-808 & 0.967 & FA/AcOH, 1.5--3.0 M & 80--120 $^\circ$C & scCO$_2$ reference route & low & low \\
    MOF-801 & 0.956 & FA/AcOH, 1.5--3.0 M & 70--110 $^\circ$C & solvent exchange + slow-drying screen & low & low \\
    PCN-222 & 0.950 & FA/AcOH, 1.5--3.0 M & 90--130 $^\circ$C & solvent exchange + slow-drying screen & low & low \\
    HKUST-1 & 0.887 & low/no acid EtOH/H$_2$O gel-aging & 40--80 $^\circ$C & EtOH exchange + slow-drying screen & low & low \\
    MIL-100(Fe) & 0.879 & acid-modulated aqueous; avoid fast crystallization & 130--170 $^\circ$C & solvent exchange + slow-drying screen & low & low \\
    ZIF-8 & 0.872 & composite/scaffold before pristine claim & 25--45 $^\circ$C & EtOH exchange + slow-drying screen & low & low \\
    Mg-MOF-74 & 0.830 & modulator screen required & 90--130 $^\circ$C & solvent exchange + slow-drying screen & medium & medium \\
    MIL-101(Cr) & 0.732 & acid-modulated aqueous; avoid fast crystallization & 160--200 $^\circ$C & solvent exchange + slow-drying screen & low & medium \\
    MOF-5 & 0.597 & composite/scaffold before pristine claim & 80--120 $^\circ$C & EtOH exchange + slow-drying screen & low & medium \\
    CAU-10 & 0.500 & modulator screen required & 130--170 $^\circ$C & solvent exchange + scCO$_2$ or very slow low-$\gamma$ dry & high & high \\
    \bottomrule
  \end{tabular}
\end{table}

\subsection{Integrated interpretation of the prototype outputs}

The prototype ranking identifies a Zr-rich high-priority group--NU-1000, UiO-67, MOF-808, MOF-801, and PCN-222--with demonstration MFP-Scores between 0.950 and 0.974. This pattern is chemically plausible because these candidates combine comparatively stable Zr-based nodes, low or mitigated capillary pressure in the heuristic Young--Laplace layer, and drying routes that avoid rapid oven drying. However, the ranking remains a software-generated prioritization order, not a claim that these materials have already been converted into monoliths.

The transition-metal and imidazolate group remains more heterogeneous. HKUST-1, MIL-100(Fe), ZIF-8, Mg-MOF-74, and MIL-101(Cr) are all ranked as high-priority demonstration candidates, but the routes differ in risk origin. Mg-MOF-74 is assigned a medium capillary-risk flag despite a high MFP-Score because of its narrower effective pore-radius descriptor. MIL-101(Cr) retains a lower demonstration score than the Zr-rich candidates, reflecting less favorable synthetic-window descriptors in the current demonstration feature table. These distinctions are useful for experimental triage because they separate a high MFP-Score from low drying risk.

The two cautionary cases are MOF-5 and CAU-10. MOF-5 is no longer treated as an almost impossible candidate; the prototype assigns it a medium demonstration score (0.597) but still flags its low framework-stability descriptor. CAU-10 is assigned the lowest demonstration score (0.500) and the only high drying-risk flag because its synthetic feature row combines high capillary pressure with an oven-drying route. These outputs support keeping both materials as controls or boundary cases in a future validation set rather than as lead experimental targets.

\subsection{Prototype-generated OPR example: UiO-67}

The accompanying software package regenerates the UiO-67 entry as a high-priority hypothesis-generating candidate (MFP-Score = 0.970; drying risk = low). The resulting OPR is deliberately expressed as a screening window rather than as a finished recipe: (i) use a Zr precursor and BPDC linker, (ii) screen formic acid or acetic acid as the modulator in the 1.5--3.0 M range, (iii) test an 80--120 $^\circ$C synthesis window, (iv) retain solvent exchange followed by slow drying, and (v) start with a small-vial gelation screen plus solvent-exchange/drying controls. This OPR should replace the earlier single-point UiO-67 illustrative protocol because it is closer to what the prototype actually generates: an experimentally testable processing window with explicit uncertainty, not a finalized synthesis route.

\section{Experimental Validation Roadmap}

The AI-generated OPRs are hypotheses that require rigorous experimental validation focusing on both the intrinsic chemistry and the critical macroscopic properties.

\subsection{Validation Metrics and Instrumentation}

The validation of a monoMOF involves a multi-modal characterization battery 
to confirm both the MOF structure and the monolithic morphology. 
In situ $\text{MOF}$ nucleation studies (e.g., on polymer fibers) 
can provide valuable real-time data.\cite{ref35}

% --- Table 2-1 ---
\begin{table}[htbp]
  \centering
  \caption{Experimental Validation Metrics (Method \& Output)}
  \label{tab:table2-1}
  \begin{tabular}{p{3cm} p{4.2cm} p{8cm}}
    \toprule
    \textbf{Category} & \textbf{Test Method} & \textbf{Key Output} \\
    \midrule
    Chemical \& Crystallographic & 
    PXRD (Powder X-ray Diffraction) & 
    Determination of crystal structure, stoichiometry, and unit cell parameters. \\

    Microstructural & 
    SEM/TEM (Scanning or Transmission Electron Microscopy) & 
    Identification of primary crystallite size, intergrowth characteristics, and overall packing density. \\

    Porosity \& Surface Area & 
    $\text{N}_2$ Adsorption (BET/BJH) & 
    Quantification of micropore volume ($V_{micro}$), mesopore volume ($V_{meso}$), and specific surface area. \\

    Composition & 
    ICP-OES (Inductively Coupled Plasma – Optical Emission Spectrometry) & 
    Determination of metal-to-linker ratio and elemental composition. \\

    Macroscopic \& Mechanical & 
    Bulk Density and Compressive Strength Tests & 
    Evaluation of volumetric density, compressive modulus ($E_{comp}$), and yield stress ($\sigma_y$). \\

    Kinetic/Rheological & 
    Rheometry & 
    Measurement of storage modulus ($\text{G}'$) and loss modulus ($\text{G}''$). \\
    \bottomrule
  \end{tabular}
\end{table}


% --- Table 2-2 ---
\begin{table}[htbp]
  \centering
  \caption{Experimental Validation Metrics (Instrumentation \& Criteria)}
  \label{tab:table2-2}
  \begin{tabular}{p{3cm} p{4.2cm} p{8cm}}
    \toprule
    \textbf{Category} & \textbf{Instrumentation} & \textbf{MonoMOF Success Criteria} \\
    \midrule
    Chemical \& Crystallographic & 
    Synchrotron source or high-resolution benchtop diffractometer & 
    Retention of the MOF framework with high crystallinity 
    and minimal amorphous background. \\

    Microstructural & 
    High-resolution TEM with EDX mapping & 
    Presence of small primary nanoparticles ($<20$ nm), 
    dense intergrowth, and clear necking between particles. \\

    Porosity \& Surface Area & 
    Gas sorption analyzer (e.g., $\text{Quantachrome}$, $\text{Micromeritics}$) & 
    Coexistence of intrinsic micropores and inter-particle mesopores, 
    together with a high BET surface area. \\

    Composition & 
    Plasma-based spectrometer & 
    Verification of correct $\text{M}:\text{L}$ stoichiometry and sample purity, 
    including nanoparticle sizing\cite{ref36} 
    and detection of residual modulator. \\

    Macroscopic \& Mechanical & 
    Universal testing machine with micro-indentation attachment & 
    High bulk density ($>0.6 \text{ g} \cdot \text{cm}^{-3}$), 
    compressive modulus ($E_{comp} > 10$ MPa), 
    and absence of macroscopic cracks. \\

    Kinetic/Rheological & 
    Stress-controlled rheometer & 
    $\text{G}'$ must exceed $\text{G}''$ in the gel state, 
    confirming a solid-like viscoelastic network prior to drying. \\
    \bottomrule
  \end{tabular}
\end{table}


\subsection{Multi-Objective Optimization and Bayesian Policy}

Successful monoMOF synthesis requires balancing three competing objectives that define a Pareto frontier:
\begin{enumerate}
    \item Porosity preservation – maximizing micropore volume $V_{micro}$.
    \item Mechanical integrity – maximizing compressive modulus $E_{comp}$.
    \item Synthesis efficiency – minimizing total synthesis time $T_{synth}$.
\end{enumerate}
A future AI-guided platform could move the system toward the optimal frontier through \textit{Bayesian Optimization (BO)}, which is ideal for expensive, noisy, and black-box experimental systems.
\[
\text{Maximize } (V_{micro}, E_{comp}, -T_{synth})
\]
subject to experimental feasibility constraints.

\subsubsection{Bayesian Optimization Implementation}

BO operates by maintaining a \textit{Gaussian Process (GP)} model of the unknown objective function $f(\mathbf{x}) \rightarrow \text{Output Metric}$ and using an \textit{Acquisition Function} $\mathcal{A}(\mathbf{x})$ to decide the next point to sample.
\begin{itemize}
    \item Gaussian Process (GP): The GP models the function by estimating both the mean $\mu(\mathbf{x})$ (the predicted outcome) and the variance $\sigma^2(\mathbf{x})$ (the uncertainty) for any point $\mathbf{x}$ in the feature space.
    \item Acquisition Function (Expected Improvement): The \textit{Expected Improvement (EI)} function could be used, which balances the trade-off between \textit{Exploitation} (sampling where the mean $\mu$ is high) and \textit{Exploration} (sampling where the variance $\sigma^2$ is high). The EI function $\mathcal{A}_{EI}(\mathbf{x})$ calculates the expected improvement over the best result found so far, $f_{best}$.
\end{itemize}
\[
\mathcal{A}_{EI}(\mathbf{x}) = E[\max(0, f(\mathbf{x}) - f_{best})]
\]
A future BO loop would iteratively suggest a new set of synthesis parameters $\mathbf{x}_{next} = \arg \max_{\mathbf{x}} \mathcal{A}_{EI}(\mathbf{x})$, the experiment would be run, the result characterized, and the GP model updated. This could enable more rapid convergence to the optimal OPR with the fewest experiments.

A long-term implementation could integrate \textit{Autonomous Experimental Execution}, in which robotic platforms carry out AI-suggested reactions under human-approved safety constraints. Real-time feedback could be provided by in-line sensors such as PXRD, which would monitor the crystalline structure of the MOF, and FTIR, which would track chemical functional groups and reaction progress. Such continuous feedback would enable ongoing model refinement, but a fully closed-loop \textit{Self-Driving Chemistry} platform is not claimed in the present work.\cite{ref19}

\subsection{Detailed Rheological and Viscoelastic Modeling of MOF Gels}

The transition from a liquid sol to a solid gel is defined by a critical change in viscoelastic properties, which is essential to model accurately for the OPR. Rheology provides the most direct measure of the gel's internal network structure and its ability to withstand stress. The complexity of these systems is analogous to that of supramolecular hydrogels\cite{ref37} and general \textit{yielding liquids}.\cite{ref38}

\subsubsection{Defining the Gel Point and Critical Exponent Scaling}

The \textit{gel point} (or critical gelation concentration, $C_{gel}$) is rigorously defined as the point where the storage modulus ($\text{G}'$) and the loss modulus ($\text{G}''$) become independent of the oscillation frequency ($\omega$). At this critical point, the complex modulus $G^*$ exhibits power-law behavior:
\[
G^*(\omega) = G_{crit} \cdot \omega^n
\]
where $n$ is the \textit{critical relaxation exponent} ($\text{tan}(\delta) = \text{G}'' / \text{G}'$) and $G_{crit}$ is the gel strength. For a near-ideal polymeric gel, $n$ is typically around $0.5$. For MOF gels, which are particle-based networks, $n$ is often lower ($0.1 < n < 0.4$), indicating a more rigid, fractal structure.

The rheology features included in the AI model are:
\begin{enumerate}
    \item Gelation Time ($\tau_{gel}$): The time required for $\text{G}' > \text{G}''$.
    \item Steady-State Moduli ($\text{G}'_{eq}, \text{G}''_{eq}$): The plateau values of the moduli after gelation.
    \item Critical Exponent ($n$): The measure of the network's fractal dimension and rigidity.
\end{enumerate}

\subsubsection{Non-Newtonian Flow Models for Industrial Pumping}

For industrial scalability, the MOF gel must be pumpable (i.e., it must flow under sufficient shear stress) into molds or continuous flow reactors before it fully sets. MOF gels are typically \textit{shear-thinning non-Newtonian fluids}, meaning their viscosity decreases under shear stress ($\dot{\gamma}$).

The \textit{Herschel–Bulkley Model} accurately describes the shear-thinning behavior with a yield stress ($\tau_0$):
\[
\tau = \tau_0 + K \dot{\gamma}^n
\]
where $\tau$ is the shear stress, $\tau_0$ is the \textit{yield stress} (the minimum stress required to initiate flow), $K$ is the consistency index, and $n$ is the flow index (with $n<1$ for shear-thinning).

The yield stress $\tau_0$ is a proposed feature for a future AI model: a high $\tau_0$ indicates a strong, well-formed gel network that can withstand the initial capillary forces $P_c$ during drying. The intended relationship is:
\[
\text{MFP-Score} \propto \frac{\tau_0}{\gamma}
\]
The OPR search maximizes this ratio by suggesting parameters that increase network strength ($\tau_0$ via aging time, $\phi$) while minimizing capillary stress ($\gamma$ via solvent exchange).

\subsubsection{Micellar and Aggregation Kinetics Modeling}

The particle-based nature of MOF gelation involves complex aggregation kinetics. Early-stage growth involves the formation of primary MOF nanoparticles, followed by the reversible aggregation of these particles into chains or clusters (flocculation), leading to the final network. This process can be modeled using the \textit{Smoluchowski kinetic model} for flocculation, where the rate of aggregation is governed by the collision frequency and the stability factor $W$ (derived from the DLVO potential, Section 2.1). The AI system uses the $V_{total}$ prediction from DLVO theory to estimate $W$ and, subsequently, the theoretical aggregation rate, providing another physics-based feature for the MFP-Score prediction.

The integration of these rheological and kinetic models provides the AI with a comprehensive, multi-scale understanding of the material's behavior, ensuring the OPRs are not just chemically feasible but also mechanically robust and industrially processable.

\subsection{\textit{In Situ} and \textit{Operando} Characterization for Model Validation}

A critical component of the validation roadmap is the use of \textit{in situ} and \textit{operando} techniques to directly observe the sol-gel transition, providing real-time data to validate the AI's kinetic and mechanistic predictions. This moves beyond \textit{ex situ} characterization of the final product to understand the \textit{process} itself.
\begin{itemize}
    \item \textit{In Situ} Small/Wide-Angle X-ray Scattering (SAXS/WAXS): Using a synchrotron source, SAXS can track the evolution of nanoparticle size, shape, and inter-particle distance ($r_p$) from the initial sol. Simultaneously, WAXS tracks the emergence of crystalline Bragg peaks, confirming the exact moment of nucleation ($J$) and the rate of crystallite growth. This provides direct experimental validation for the nucleation and aggregation models ($k_{agg}$, $k_{diss}$) and the AI's predicted $\tau_g$. This data is essential for validating the \textit{latent variables} in the causal model (DAG), such as "Crystallite Size."
    \item \textit{In Situ} Spectroscopy (Raman, FTIR): By placing probes directly in the reaction vessel, these techniques monitor the chemical environment in real-time. FTIR can track the deprotonation of the linker's carboxylic acid groups and the formation of metal-carboxylate coordination bonds. Raman spectroscopy is highly sensitive to the vibrational modes of the metal-oxide clusters (e.g., the $\text{Zr}_6$ SBU), allowing for the tracking of cluster formation and how it is affected by modulators ($\alpha$) and temperature ($T$). This provides a direct ground truth for the TST-predicted kinetic constants ($k_f$).\cite{ref53}
    \item \textit{Operando} Rheo-SAXS: This powerful hyphenated technique combines rheometry (Section 5.3) with SAXS, allowing for simultaneous measurement of the evolving viscoelastic properties ($\text{G}', \text{G}''$) and the nanoscale structure. This directly links the macroscopic gel strength ($\tau_0$) to the underlying nanoparticle network formation (fractal dimension, particle size), providing the richest possible dataset for validating the MFP-Score's underlying assumptions, particularly the $\text{MFP-Score} \propto \tau_0 / \gamma$ relationship.
\end{itemize}
In a future experimental deployment, this real-time data could be fed back into the AI model in the Active Learning loop (Section 3.3), allowing the model to update its kinetic parameters ($\tau_g$, $k_f$) and move toward a more physically accurate and experimentally testable OPR.

\section{Discussion and Broader Implications}

The proposed AI-accelerated monoMOF framework aims to shift discovery from serendipity toward structured, testable prediction. LLMs can act as literature-organization and hypothesis-generation tools, helping uncover possible patterns and analogous synthesis conditions across disparate MOF families. This shift to data-defined chemistry is fundamental.

\subsection{AI-Driven Insights and Challenges}

The integration of physical chemistry principles into the ML framework (via $P_c$ and $\tau_g$) is intended to ensure that future predictions are not only data-driven but also physically interpretable. The complexity of the sol–gel process—involving non-equilibrium thermodynamics and complex particle interactions—means that any deployed ML model must capture nuances missed by classical models. The GBT approach is well suited for handling heterogeneous feature spaces containing mixtures of continuous, categorical, and high-dimensional vector features, but its performance must be established by future curated data and validation beyond the current synthetic-curated prototype before any validated accuracy claim is made.

However, challenges include the small size of the positive monoMOF dataset, necessitating unsupervised learning and \textit{transfer learning} from related material systems (e.g., silica or titania sol–gel chemistry). Future work could integrate the ML model with physical simulations (DFT, Molecular Dynamics (MD)) to refine parameter estimation and more accurately predict mechanical properties. Specifically, future MD simulations could model the inter-particle necking process during the gel-aging step, providing a \textit{Microstructural Cohesion Index} to further enhance the MFP-Score.

\subsection{Industrial Prospects: Scalability, Economics, and Sustainability}

The ability to predictably generate monoMOFs carries broad implications for industrial materials science, shifting the focus from lab-scale novelty to large-scale, process-optimized manufacturing.\cite{ref39}

\subsubsection{Cost-Benefit Analysis and Economic Viability}

The primary cost reduction comes from the acceleration of R\&D. In a future validated implementation, Bayesian optimization could reduce the number of empirical optimization experiments required to discover a robust OPR, thereby shortening the \textit{Time-to-Market} for a new MOF product. In the present manuscript, OPRs are treated as hypothesis-generating route suggestions that can be screened for scalability-relevant factors, including manufacturing and shaping processes:\cite{ref40}
\begin{itemize}
    \item Solvent Selection: A future validated implementation could prioritize \textit{Green Solvents} (water, ethanol) where possible, or suggests OPRs that minimize the volume of high-cost/high-toxicity solvents (e.g., $\text{DMF}, \text{DEF}$) through multiple solvent exchange steps.
    \item Modulator/Precursor Cost: A future validated OPR search could target the minimum effective concentration of expensive or corrosive modulators ($\text{HF}, \text{HCl}$) while maintaining the desired kinetic control, directly lowering the raw material costs.
    \item Energy Efficiency: By optimizing the synthesis temperature ($T$) and time ($T_{synth}$), a future validated OPR search could minimize the energy footprint of the reaction (e.g., preferring $120^\circ \text{C}$ for 24 hours over $180^\circ \text{C}$ for 72 hours if the OPR is stable).
\end{itemize}
The potential industrial value would lie in the final product's performance: the high bulk density and low pressure drop of monoMOFs translate directly into lower operational costs and higher space-time yields in industrial reactors compared to powder MOFs. The mechanical robustness ($E_{comp}$) could allow for simpler reactor designs without the need for sophisticated filtration or particle retention systems.

\subsubsection{Green Chemistry and Sustainability}

The AI framework is explicitly designed to incorporate sustainability constraints, guiding the search toward \textit{Green Chemistry OPRs}:
\begin{enumerate}
    \item Atom Economy Maximization: Prioritizing precursors that maximize the incorporation of all reactants into the final MOF structure.
    \item Waste Minimization: The Active Learning policy can be modified to include a penalty term for non-recyclable solvent use or high volumes of toxic byproducts.
    \item Safer Synthesis: The OPR generation includes a check against a safety database (e.g., NFPA rating, flash point) for all suggested solvents and modulators, suggesting safer alternatives when possible.
\end{enumerate}
For example, when predicting an OPR for $\text{UiO-66}$, the AI would strongly favor a water-based synthesis route\cite{ref41} over the traditional $\text{DMF}$ route, provided the $\text{MFP-Score}$ for the water-based route remains above the critical threshold ($\sim 0.7$). This aligns with major industrial goals, such as Yaghi's scalable, green synthesis of MOF-303 for water harvesting.\cite{ref42}

\subsubsection{Intellectual Property and Data Ownership}

A possible future intellectual-property focus would not be the MOF structures themselves, but reproducible, scalable OPRs required to manufacture them. Such OPRs could represent trade-secret process knowledge if experimentally validated. Furthermore, the intellectual property of a future AI model (the $\mathbf{w}$ weights and the GNN architecture) and the curated, structured database (DigiMOF-Mono)\cite{ref16} derived from LLM text mining\cite{ref3} could become valuable assets. Licensing a validated predictive model for new MOF structures could become a significant business model, offering a "Monolith-ability Score (MFP-Score) as a Service" to materials manufacturers.

\subsubsection{Applications Landscape Expansion}

The predictive power of the AI framework accelerates the realization of applications beyond traditional MOF domains:\cite{ref43}
\begin{itemize}
    \item Catalysis: Monolithic architectures permit efficient mass transfer and pressure drop minimization, enabling practical catalytic reactor integration. The hierarchical porosity allows for tandem catalysis where the micropores perform the chemical reaction and the mesopores facilitate the rapid transport of large molecular substrates. For example, a monolithic UiO-66-NH$_2$ could be used in a flow reactor for the Knoevenagel condensation, exhibiting higher stability and throughput than a packed powder bed.
    \item Gas Storage and Separation: Increased bulk density and preserved microporosity lead to record volumetric uptakes for $\text{CH}_4$ and $\text{CO}_2$, vital for energy applications.\cite{ref4} MOF-74-family adsorbents, including amine-appended Mg-MOF-74 analogues, provide well-established chemistry for $\text{CO}_2$ capture, while their translation into monolithic or shaped bodies should be treated as an engineering target rather than as a completed demonstration in this manuscript.\cite{ref44} Hydrogen storage likewise remains a long-term application driver for densified or shaped MOF architectures, where volumetric rather than purely gravimetric metrics become central.\cite{ref45}
    \item Optical and Electronic Devices: Transparent or semi-transparent monoMOFs can host functional dyes or nonlinear optical materials for direct device integration into micro-electronic components, such as MOF-based sensors deposited directly onto silicon wafers without binder interference.
    \item Chromatography: Binder-free, macroporous monoliths offer superior separation efficiency due to reduced resistance to mass transfer compared to traditional packed columns. The AI's ability to precisely tune the inter-particle pore size ($r_p$) is critical for optimizing chromatographic resolution for separating large molecules like proteins or peptides.
    \item Biomedical Implants: The mechanical robustness of monoMOFs opens the door for their use as solid, implantable devices for controlled drug delivery. A monolithic MOF loaded with a therapeutic could be implanted \textit{in vivo}, releasing the drug at a zero-order rate determined by the framework's degradation or the drug's diffusion through the hierarchical pore network.
    \item Environmental Remediation: MonoMOFs can function as `sponges' or flow-through filters in water purification systems, efficiently capturing heavy metals and persistent organic pollutants such as per- and polyfluoroalkyl substances (PFAS), a class of highly stable, bioaccumulative compounds of significant environmental concern. The monolithic architecture prevents the leaching of nanoparticles into the environment and facilitates straightforward regeneration and reuse, addressing a major limitation associated with conventional powder-based sorbents.

\end{itemize}
The shift in focus from powder to monolithic morphology, supported by auditable AI-guided prioritization, could help unlock industrially relevant MOF formats if validated experimentally.

\section{Future Manufacturing Translation: Continuous Flow Synthesis for OPR Implementation}

A future transition from AI-suggested OPRs to industrial production would require adapting reaction protocols from lab-scale batch autoclaves to \textit{continuous flow chemistry}. Flow reactors offer control over mixing, heat transfer, and residence time, which makes them an appropriate future manufacturing layer for the kinetic variables considered by the framework.\cite{ref39}

\subsection{Principles of Continuous Flow Synthesis for MOFs}

Batch synthesis typically relies on slow heating and cooling in a large volume, leading to uncontrolled concentration gradients. Continuous flow reactors, such as \textit{microfluidic reactors} or \textit{continuous stirred-tank reactor (CSTR) cascades}, allow for highly controlled, steady-state synthesis.

Critical Flow Parameters:
\begin{enumerate}
    \item Residence Time ($\tau_{res}$): Directly replaces batch reaction time ($T_{rxn}$). In a plug-flow reactor (PFR), $\tau_{res} = V_{reactor} / Q$, where $Q$ is the volumetric flow rate. The AI-derived $T_{rxn}$ is directly translated to $\tau_{res}$.
    \item Mixing Quality (Reynolds Number, $\text{Re}$): Low $\text{Re}$ (laminar flow) in microfluidics guarantees uniform mixing at the molecular level, critical for preventing localized concentration spikes that cause rapid, heterogeneous nucleation (powder formation).
    \item Temperature Homogeneity: High surface-area-to-volume ratio in flow reactors ensures instantaneous temperature control, eliminating the temperature gradients that plague large batch autoclaves.
\end{enumerate}

\subsection{Microfluidic Reactor Design for Nanoparticle Nucleation}

The initial nucleation phase, governed by the rate constant $k_f$ (Section 2.5.1), is most sensitive to concentration. Microfluidics are ideal for this stage:
\begin{itemize}
    \item Y-Junction Mixing: Reactants (Metal precursor, Linker/Modulator solution) are mixed at a Y-junction. The short diffusion distance ensures rapid, homogeneous mixing, triggering instantaneous nucleation.
    \item Growth Channel: The mixed solution flows through a heated channel (residence time 1 to 10 minutes) to control the initial growth phase, ensuring the formation of monodisperse nanocrystals ($<15 \text{ nm}$) as required for optimal gel formation.
\end{itemize}

\subsection{Continuous Gelation and Extrusion}

The subsequent gelation and aging phase requires controlled, slow mixing to promote inter-particle necking ($\text{OA}$ and $\text{PMC}$).
\begin{itemize}
    \item CSTR Cascade: The solution is transferred from the microfluidic unit to a cascade of 2-3 CSTRs, where the agitation rate is precisely controlled to maintain the particles just above the critical aggregation concentration ($CAC$). The total residence time in the cascade implements the AI-suggested gel aging time ($\tau_{gel}$).
    \item Viscosity Monitoring: In-line rheometers constantly monitor the gel's viscoelastic state ($\text{G}'/\text{G}''$ ratio) at the cascade outlet. This real-time data feeds directly back into the RL/BO loop to adjust $Q$ or $T$, maintaining the optimal gel point before the next processing step.
\end{itemize}
Continuous Molding and Drying:
In a future scale-out scenario, a homogeneous gel could be continuously pumped ($\tau_0$ requirement met) through an extrusion nozzle into molds or structured beds. Final solvent exchange and controlled drying (e.g., using a temperature/humidity controlled tunnel oven to mimic the low $R_d$ requirement) would then be treated as reactor-specific engineering variables. This end-to-end concept is a manufacturing translation target rather than a completed industrial process in the present manuscript.

\subsection{Addressing Heat and Mass Transfer in Scale-Up}

Scaling up from laboratory-scale (milliliters) to pilot-plant-scale (liters) frequently introduces significant heat and mass transfer limitations that compromise the assumptions underlying kinetic models.

\begin{itemize}
    \item Mass transfer: In large batch reactors, the diffusion of linkers to the growing crystal surfaces often becomes rate-limiting, resulting in concentration depletion zones and increased crystal size polydispersity in powders. In contrast, flow reactors mitigate this issue through enhanced surface area and convective transport. In a future reactor-specific model, the diffusion term, $D \nabla^2 C$ (Section 2.2), would be adjusted to account for the effective enhancement in mixing under flow conditions ($D_{\text{eff}} \gg D_{\text{molecular}}$).\cite{ref39}
    \item Heat transfer: The strongly exothermic nature of MOF crystallization can generate hotspots in large batch reactors, leading to localized thermal degradation or uncontrolled nucleation. Flow reactors facilitate rapid heat dissipation via microchannel walls, maintaining near-isothermal conditions and preserving the AI-predicted temperature $T$ with high fidelity. This thermal stability is essential for the accurate application of the TST-derived kinetic constant $k_f$.
\end{itemize}


\subsection{Process Analytical Technology (PAT) Integration}

To move toward a future self-driving laboratory, a continuous-flow reactor would need to be integrated with \textit{Process Analytical Technology (PAT)}, as required by regulatory standards for modern manufacturing. Such integration would leverage the \textit{in situ} sensors discussed in Section 5.4 (Raman, FTIR, SAXS) to provide real-time process feedback.

\begin{itemize}
    \item PAT feedback loop: Should the \textit{in situ} Raman probe detect a deviation in the modulator's protonation state—a leading indicator of altered reaction kinetics—the AI system can immediately adjust the modulator flow rate. Similarly, if \textit{in situ} SAXS measurements reveal particle size drift (e.g., particles measure 20~nm rather than the target 15~nm), the AI can modulate the reactor temperature or residence time ($Q$) to restore the process to the OPR setpoint. This PAT-driven control represents the practical implementation of a Digital Twin (Section 8.3) and is critical for maintaining product quality and consistency at an industrial scale.
\end{itemize}

By explicitly embedding these chemical-engineering principles, a future validated AI framework could help screen whether predicted OPRs are chemically plausible and potentially scalable, while final industrial deployment would require reactor-specific experiments and process validation.


\section{Future Directions: Multi-Scale Modeling and Hybrid AI}

To push the frontiers of monoMOF discovery, the AI framework must evolve toward true multi-scale modeling and hybrid generative AI structures. This work is a crucial part of the broader efforts in the future of computational chemistry.\cite{ref46}

\subsection{Generative AI for Linker and Node Design}

Current AI efforts in this manuscript primarily focus on organizing and prioritizing the \textit{process} ($\mathbf{x}_{OPR}$) for \textit{existing} MOFs. A future frontier would involve leveraging generative models to propose \textit{entirely new linkers} ($\mathbf{L}_{new}$) and \textit{nodes} ($\mathbf{M}_{new}$) with high predicted MFP-Score priors. Data-driven predictions of MOF structures provide a robust foundation for this generative design strategy.\cite{ref47}

\begin{itemize}
    \item Constrained generation: The generative model is guided by features associated with high MFP-Scores. Linkers are required to be rigid (high GNN flexibility index), exhibit strong binding groups (high ESP at coordination sites), and possess a predicted $E_{gap}$ within a thermodynamically stable range. This constitutes an \textit{inverse design} problem: the desired \textit{properties} (e.g., $E_{comp} > 20~\text{MPa}$, $\text{MFP-Score} > 0.9$) are specified, and the AI generates the corresponding \textit{materials} that satisfy these criteria.\cite{ref48,ref49} Typical implementations utilize generative architectures such as \textit{Variational Autoencoders (VAEs)} or \textit{Generative Adversarial Networks (GANs)} operating on SMILES strings or graph representations. More advanced 3D-aware generative models, including diffusion- or flow-based frameworks, can directly generate linker geometries pre-optimized for specific framework topologies.\cite{ref54}
    \item Node substitution: The model can propose novel metal nodes (e.g., $\text{Hf}_6$, $\text{V}_4$) exhibiting high $\Delta E_{diss}$ (Node Stability Index) and favorable coordination geometries for dense packing, thereby promoting robust intergrowth. Further constraints, including cost and toxicity considerations, guide the AI to generate linkers and nodes that are not only structurally effective but also economically feasible and environmentally benign.
\end{itemize}

In a future implementation, this extension would move the AI system from process optimization toward hypothesis generation for candidate linkers and framework variants based on target macroscopic performance metrics ($E_{comp}, V_{meso}$) and process constraints such as the use of green solvents.


\subsection{Integrating Ab Initio Molecular Dynamics (AIMD)}

While classical MD (Section 2.5.2) provides elastic constants, it relies on parameterized force fields. In contrast, \textit{Ab Initio Molecular Dynamics (AIMD)}, which calculates forces directly from the electronic structure via DFT at each time step, delivers the highest-fidelity simulations.

\begin{itemize}
    \item AIMD for modulator effects: AIMD explicitly captures the dynamic coordination environment of the metal cluster in the presence of modulators (e.g., $\text{HF}$ or $\text{Formic Acid}$) and solvent molecules. This enables rigorous, temperature-dependent evaluation of instantaneous cluster size and stability, thereby refining the Node Stability Index beyond static DFT optimizations.
    \item Proton transfer dynamics: AIMD can model proton transfer between modulators and linkers, representing the initial kinetic step of the reaction. This provides a quantum-accurate, dynamic estimate of the activation barrier, $\Delta G^\ddagger$, which surpasses the static TST approximation.
\end{itemize}

Although AIMD is computationally prohibitive for high-throughput screening, it remains the gold standard for validating the most promising OPRs and for training next-generation, lower-cost, high-accuracy ML force fields.

\subsection{Digital Twin of the Synthesis Reactor}

The ultimate objective is the creation of a \textit{Digital Twin}—a real-time, high-fidelity virtual representation of the continuous flow reactor (Section 7).

\begin{itemize}
    \item Fluid dynamics (CFD): Integration of the reaction–diffusion equation (Section 2.2) with \textit{Computational Fluid Dynamics (CFD)} simulations provides a spatially and temporally resolved map of concentration, temperature, and shear stress throughout the reactor geometry.
    \item Future real-time prediction: In a later implementation, rheological and in-line spectroscopic data from a physical reactor (via PAT, Section 7.5) could be used by a Digital Twin to estimate whether gelation is likely to result in a monolith and to recommend control adjustments (e.g., modulating flow rate $Q$ or temperature $T$). This would operationalize the SCM framework (Section 3.4) as a future in-silico decision-support layer rather than as a completed capability in the present prototype.
\end{itemize}

By integrating LLMs, graph-based representations, causal inference, optional quantum simulations, and chemical engineering principles, this approach outlines a route toward more transparent and experimentally testable monoMOF discovery rather than claiming a completed autonomous synthesis platform.



% --- APPENDICES ---
\appendix




\section*{Data and code availability}
The manuscript is accompanied by Supplementary Software~1, a reproducible software-resource package for reviewer inspection, manuscript-table regeneration, and final repository deposition. The package contains synthetic-curated demonstration data, a feature dictionary, a synthesis-extraction schema, baseline MFP-Score and OPR-ranking scripts, a capillary-stress risk calculator, generated outputs, a notebook, README and reproducibility documentation, and reviewer-facing quick-start material. For Digital Discovery-style reproducibility, the package also includes one-command regeneration and validation scripts, environment files, output manifests, checksums, ablation and sensitivity summaries, an output-validation report, a clean-run reproducibility audit, and a manuscript-output traceability map. To reduce reviewer friction, the submission package also contains root-level package-visibility files, including a package contents summary, a reviewer ZIP readme, and a software file index.

The package is designed so that a reviewer can regenerate the manuscript-facing MFP-score ranking, OPR table, feature-diagnostic table, uncertainty report, ablation summaries, and sensitivity summaries from the supplied demonstration data. The code-derived outputs remain demonstration-level and hypothesis-generating: they do not constitute a validated production model, real DFT/MD workflow, autonomous-lab implementation, or experimentally verified synthesis route. Final quantitative performance claims, final candidate rankings, or final code-derived feature-importance values should be regenerated from a curated dataset before being presented as predictive results.

For final submission, the exact software release, data tables, environment file, output manifests, checksums, reviewer quick-start documentation, package-visibility files, graphical-abstract/TOC brief, and software-resource documentation should be archived in a persistent repository such as Zenodo or an equivalent institutional repository. The release package is prepared as a DOI-ready bundle with Zenodo/GitHub metadata drafts, a reviewer-facing manifest, and confirmed MIT License files at both the repository root and the Supplementary Software 1 level. Final repository DOI placeholder: [INSERT VERSIONED ZENODO DOI FOR FINAL RELEASE]. Final software license line: MIT License. For the present RSC-compatible submission package, the scientific content and software outputs are unchanged from the validated software baseline.

\section*{Acknowledgments}
\setlength{\parindent}{1em}  % 首行缩进
\setlength{\parskip}{0.5em}  % 段间距
The author acknowledges collaboration discussions with Prof. Andrew Wheatley (University of Cambridge), whose work on MOF synthesis and AI-assisted chemistry inspired this framework. No production-scale DFT, MD, autonomous-lab, or deployed software computation is claimed in the present manuscript.

\section*{Conflict of Interest}
\setlength{\parindent}{1em}
\setlength{\parskip}{0.5em}
The authors declare no competing financial interests.

\section*{Author}
\setlength{\parindent}{1em}
\setlength{\parskip}{0.5em}
Jinnan Wei, Department of Chemistry, University of Cambridge, UK.

\section*{Keywords}
\setlength{\parindent}{1em}
\setlength{\parskip}{0.5em}
MOFs, monolithic frameworks, AI-guided synthesis, machine learning, sol--gel process, LLM, materials informatics, MFP-Score, Optimal Processing Route, reproducibility, Green Chemistry, Bayesian Optimization, Graph Neural Networks, Causal Inference, DFT roadmap, Molecular Dynamics roadmap, Continuous Flow Synthesis, Rheology.


% --- REFERENCES ---
% RSC-style numbered references are kept in a manual thebibliography block for portability.
% Final production can convert this block to BibTeX if requested by the journal portal.

\begin{thebibliography}{99}

\bibitem{ref1} Furukawa, H.; Cordova, K. E.; O'Keeffe, M.; Yaghi, O. M. The Chemistry and Applications of Metal--Organic Frameworks. \textit{Science} \textbf{2013}, \textit{341}, 1230444. DOI: 10.1126/science.1230444.

\bibitem{ref2} Tian, T.; Zhang, J.; Jiang, Z.; Chen, Z.; Wheatley, A. E. H.; Wang, B. Synthesis and Applications of Monolithic Metal-Organic Frameworks. \textit{Chem. Soc. Rev.} \textbf{2021}, \textit{50}, 12767–12816.

\bibitem{ref3} Zheng, Z.; Zhang, O.; Borgs, C.; Chayes, J. T.; Yaghi, O. M. ChatGPT Chemistry Assistant for Text Mining and Prediction of MOF Synthesis. \textit{arXiv} \textbf{2023}, arXiv:2306.11296.

\bibitem{ref4} Connolly, B. M.; Madden, D. G.; Wheatley, A. E. H.; Fair, A.; Scott, M. J.; O’Brien, J. E.; Gascón, J.; Moghadam, P. Z.; Souto, M.; Clowes, R.; et al. Tuning Porosity in Macroscopic Monolithic Metal-Organic Frameworks for Exceptional Natural Gas Storage. \textit{Nat. Commun.} \textbf{2019}, \textit{10}, 2364.

\bibitem{ref5} Madden, D. G.; Connolly, B. M.; Clowes, R.; Gascón, J.; O’Brien, J. E.; Wheatley, A. E. H.; Moghadam, P. Z.; Fairen-Jimenez, D. Monolithic Metal-Organic Framework Materials for Gas Separations and Carbon Capture. \textit{Faraday Discuss.} \textbf{2021}, \textit{231}, 246–260.

\bibitem{ref6} Geng, H.; Guan, X.; Lu, C.; Li, J. Metal–Organic Framework Gels and Monoliths. \textit{Chem. Sci.} \textbf{2020}, \textit{11}, 837–846.

\bibitem{ref7} Wheatley, A. E. H.; Connolly, B. M.; Johnson, E. L.; O'Brien, J. E.; Fairen-Jimenez, D. In Situ Immobilization of SnO$_2$ Nanoparticles in a Monolithic ZIF-8 (monoZIF-8) Support. \textit{Adv. Funct. Mater.} \textbf{2018}, \textit{28}, 1803276.


\bibitem{ref9} Liu, B.; Li, S.; Wang, R.; Zhang, S.; Li, S. Recent Advances in the Synthesis of Monolithic Metal–Organic Frameworks. \textit{J. Mater. Chem. A} \textbf{2021}, \textit{9}, 5947–5963.

\bibitem{ref10} Park, H.; Kang, Y.; Choe, W.; Kim, J. Mining Insights on Metal--Organic Framework Synthesis from Scientific Literature Texts. \textit{arXiv} \textbf{2021}, arXiv:2108.13590.

\bibitem{ref11} Brinker, C. J.; Scherer, G. W. \textit{Sol--Gel Science: The Physics and Chemistry of Sol--Gel Processing}; Academic Press: San Diego, 1990.

\bibitem{ref12} Pierre, A. C. \textit{Introduction to Sol--Gel Processing}; Springer: New York, 1998.

\bibitem{ref13} Moosavi, S. M.; et al. From Data to Discovery: Recent Trends of Machine Learning in Metal–Organic Frameworks. \textit{JACS Au} \textbf{2024}, \textit{4}, 3326–3342.

\bibitem{ref14} Jablonka, K. M.; Moosavi, S. M.; Asgari, M.; Smit, B. Machine Learning in Materials Science: A Case Study of Metal–Organic Frameworks. \textit{MRS Bull.} \textbf{2021}, \textit{46}, 558–566.

\bibitem{ref15} Davis, M. E.; Lobo, R. F. Zeolite and Zeotype Materials: A Historical Perspective and Future Outlook. \textit{Chem. Mater.} \textbf{2020}, \textit{32}, 1515–1527.

\bibitem{ref16} Glasby, L. T.; Gubsch, K.; Bence, R.; Oktavian, R.; Isoko, K.; Cordiner, J. L. DigiMOF: A Database of Metal--Organic Framework Synthesis Information Generated via Text Mining. \textit{Chem. Mater.} \textbf{2023}, \textit{35}, 5099--5109.

\bibitem{ref17} Kalidindi, S. R.; Medford, A. J. Predicting and Accelerating Nanomaterials Synthesis Using Machine Learning Featurization. \textit{Annu. Rev. Mater. Res.} \textbf{2020}, \textit{50}, 367–391.

\bibitem{ref18} Jha, D.; Ward, L.; Paul, A.; Liao, W.-k.; Choudhary, A.; Wolverton, C.; Agrawal, A. ElemNet: A Deep Learning Framework for Materials Discovery. \textit{Sci. Rep.} \textbf{2019}, \textit{9}, 17593.

\bibitem{ref19} Burger, B.; Maffettone, P. M.; Gusev, V. V.; Aitchison, C. M.; Bai, Y.; Wang, X.; Li, X.; Alston, B. M.; Li, B.; Clowes, R.; et al. A Mobile Robotic Chemist. \textit{Nature} \textbf{2020}, \textit{583}, 237–241.

\bibitem{ref20} Pearl, J. \textit{Causality: Models, Reasoning, and Inference}; Cambridge University Press, \textbf{2009}.

\bibitem{ref21} Schütt, K. T.; Gastegger, M.; Tkatchenko, A.; Müller, K.-R.; Gasteiger, J. Unifying Machine Learning and Quantum Chemistry with a Deep Neural Network for Molecular Wavefunctions. \textit{Nat. Commun.} \textbf{2019}, \textit{10}, 5024.

\bibitem{ref22} Klicpera, J.; Groß, J.; Günnemann, S. Directional Message Passing for Molecular Graphs. \textit{arXiv} \textbf{2020}, arXiv:2003.03123.

\bibitem{ref35} Mandemaker, L. D. B.; Filez, M.; Delen, G.; Tan, H.; Zhang, X.; Lohse, D.; Weckhuysen, B. M. Time-Resolved In Situ Liquid-Phase Atomic Force Microscopy and Infrared Nanospectroscopy during the Formation of Metal--Organic Framework Thin Films. \textit{arXiv} \textbf{2020}, arXiv:2002.00683.

\bibitem{ref36} Prause, A.; et al. Validation of Single Particle ICP-MS for Routine Measurements of Nanoparticle Size and Number Size Distribution. \textit{Anal. Chem.} \textbf{2019}, \textit{91}, 1634–1641.

\bibitem{ref37} Yan, C.; Pochan, D. J. Rheological Properties of Peptide-Based Hydrogels for Biomedical and Other Applications. \textit{Chem. Soc. Rev.} \textbf{2010}, \textit{39}, 3528--3540.

\bibitem{ref38} Bonn, D.; Denn, M. M.; Berthier, L.; Divoux, T.; Manneville, S. Yield Stress Materials in Soft Condensed Matter. \textit{Rev. Mod. Phys.} \textbf{2017}, \textit{89}, 035005.

\bibitem{ref39} James, S. L.; et al. Continuous Flow Synthesis of Metal-Organic Frameworks: A Review of the Last 5 Years. \textit{Chem. Soc. Rev.} \textbf{2019}, \textit{48}, 4423–4448.

\bibitem{ref40} Wang, B.; et al. Binding Materials for MOF Monolith Shaping Processes: A Review towards Real Life Application. \textit{Energies} \textbf{2022}, \textit{15}, 1489.

\bibitem{ref41} Van de Voorde, B.; et al. A Water-Based and Modulator-Free Synthesis of UiO-66 and Its Application in Gas-Phase Separations. \textit{Green Chem.} \textbf{2015}, \textit{17}, 3184–3188.

\bibitem{ref42} Hanikel, N.; et al. Scalable and Green Synthesis of the Water-Harvester MOF-303. \textit{ACS Sustainable Chem. Eng.} \textbf{2021}, \textit{9}, 9010–9017.

\bibitem{ref43} Li, H.; et al. Beyond Powders: Monolithic Devices on the Basis of Metal-Organic Frameworks (MOFs). \textit{Angew. Chem. Int. Ed.} \textbf{2021}, \textit{60}, 1070–1086.

\bibitem{ref44} McDonald, T. M.; Mason, J. A.; Kong, X.; Bloch, E. D.; Gygi, D.; Dani, A.; Crocella, V.; Giordanino, F.; Odoh, S. O.; Drisdell, W. S.; Vlaisavljevich, B.; Dzubak, A. L.; Poloni, R.; Schnell, S. K.; Planas, N.; Lee, K.; Pascal, T.; Wan, L. F.; Prendergast, D.; Neaton, J. B.; Smit, B.; Kortright, J. B.; Gagliardi, L.; Bordiga, S.; Reimer, J. A.; Long, J. R. Cooperative Insertion of CO2 in Diamine-Appended Metal--Organic Frameworks. \textit{Nature} \textbf{2015}, \textit{519}, 303--308.

\bibitem{ref45} Murray, L. J.; Dinca, M.; Long, J. R. Hydrogen Storage in Metal--Organic Frameworks. \textit{Chem. Soc. Rev.} \textbf{2009}, \textit{38}, 1294--1314.

\bibitem{ref46} Deringer, V. L.; et al. Machine Learning and the Future of Computational Chemistry. \textit{Nat. Rev. Chem.} \textbf{2019}, \textit{3}, 287–296.

\bibitem{ref47} Li, Z.; et al. Data-Driven Prediction of Structures of Metal-Organic Frameworks. \textit{J. Am. Chem. Soc.} \textbf{2024}, \textit{146}, 9789–9798.

\bibitem{ref48} Kim, J.; et al. Inverse Design of Metal–Organic Frameworks Using a Deep Generative Model. \textit{Nat. Commun.} \textbf{2020}, \textit{11}, 5311.

\bibitem{ref49} Sánchez-Lengeling, B.; Aspuru-Guzik, A. Inverse Molecular Design Using Machine Learning: Generative Models for De Novo Drug and Materials Discovery. \textit{Acc. Chem. Res.} \textbf{2018}, \textit{51}, 1243–1252.

\bibitem{ref50} De Yoreo, J. J.; Gilbert, P. U. P. A.; Sommerdijk, N. A. J. M.; Penn, R. L.; Whitelam, S.; Joester, D.; Zhang, H.; Rimer, J. D.; Navrotsky, A.; Banfield, J. F. Crystallization by Particle Attachment. \textit{Science} \textbf{2015}, \textit{349}, aaa6760.

\bibitem{ref51} Addicoat, M. A.; Vankova, N.; Akter, I. F.; Heine, T. Extension of the Universal Force Field to Metal--Organic Frameworks. \textit{J. Chem. Theory Comput.} \textbf{2014}, \textit{10}, 880--891.

\bibitem{ref52} Li, Y.; et al. A Coarse-Grained Molecular Dynamics Study on the Self-Assembly of Metal-Organic Framework Nanoparticles. \textit{J. Phys. Chem. C} \textbf{2019}, \textit{123}, 18605–18613.

\bibitem{ref53} Taddei, M. \textit{In Situ} and \textit{Operando} Characterization of Metal-Organic Frameworks. \textit{Chem. Soc. Rev.} \textbf{2017}, \textit{46}, 3128–3154.

\bibitem{ref54} Gebauer, N. W. A.; Gastegger, M.; Schütt, K. T. A Review on Inverse Design for Molecular and Materials Sciences. \textit{Nat. Rev. Chem.} \textbf{2023}, \textit{7}, 312–326.

\end{thebibliography}

\end{document}

